\ticket{21}{Понятие языковой модели, виды моделей. Статистическая языковая модель: униграммные и N-граммные модели, применение в задачах обработки текстов. Сглаживание модели. Перплексия и способ ее подсчета}

\begin{examplan}
    \item Дать понятие языковой модели и перечислить виды.
    \item Объяснить униграммные и n-граммные статистические модели.
    \item Описать проблему разреженности и сглаживание.
    \item Объяснить перплексию и способ ее подсчета.
\end{examplan}

\subsection*{Короткая версия для письменного ответа}

\textbf{Языковая модель} задает вероятностное распределение на последовательностях слов и позволяет оценивать $P(w_1,\ldots,w_m)$ или вероятность следующего слова по контексту.

По цепному правилу:
\[
P(w_1,\ldots,w_m)=\prod_{i=1}^{m}P(w_i\mid w_1,\ldots,w_{i-1}).
\]

\subsection*{Виды языковых моделей}

\begin{itemize}
    \item \textbf{Статистические}: основаны на частотах слов и n-грамм в корпусах.
    \item \textbf{Нейросетевые}: используют нейронные сети, эмбеддинги и контекстные представления.
    \item \textbf{Авторегрессионные}: предсказывают следующее слово по предыдущим.
    \item \textbf{Маскированные}: предсказывают скрытые токены по левому и правому контексту.
\end{itemize}

\subsection*{N-граммные модели}

N-грамма --- последовательность из $N$ слов. N-граммная модель делает марковское предположение: вероятность слова зависит только от предыдущих $N-1$ слов:
\[
P(w_i\mid w_1,\ldots,w_{i-1})\approx P(w_i\mid w_{i-N+1},\ldots,w_{i-1}).
\]

Оценка максимального правдоподобия:
\[
P(w_i\mid context)=\frac{Count(context,w_i)}{Count(context)}.
\]

\begin{description}
    \item[Униграммная модель] игнорирует контекст: $P(W)=\prod_i P(w_i)$.
    \item[Биграммная модель] учитывает одно предыдущее слово.
    \item[Триграммная модель] учитывает два предыдущих слова.
\end{description}

Чем больше $N$, тем больше контекст, но тем сильнее проблема разреженности: многие допустимые n-граммы не встречаются в обучающем корпусе.

\subsection*{Применение}

N-граммные модели применяются в распознавании речи, машинном переводе, автодополнении, проверке орфографии, информационном поиске и как признаки в классификации текстов.

\subsection*{Сглаживание}

\textbf{Сглаживание} перераспределяет вероятностную массу от частых n-грамм к редким и невстречавшимся, чтобы избежать нулевых вероятностей.

Методы:
\begin{itemize}
    \item \textbf{Лапласовское сглаживание}: добавить 1 к счетчикам;
    \item \textbf{add-$\delta$}: добавить малое число $\delta$;
    \item \textbf{интерполяция}: смешать модели разных порядков;
    \item \textbf{backoff}: при ненадежной n-грамме перейти к модели меньшего порядка.
\end{itemize}

\subsection*{Перплексия}

\textbf{Перплексия} --- внутренняя метрика качества языковой модели на тестовом корпусе. Чем ниже перплексия, тем лучше модель предсказывает текст.

Для последовательности $W=(w_1,\ldots,w_M)$:
\[
PP(W)=P(W)^{-1/M}.
\]

В логарифмической форме:
\[
PP(W)=2^{-\frac{1}{M}\sum_{i=1}^{M}\log_2 P(w_i\mid w_1,\ldots,w_{i-1})}.
\]

Интуитивно перплексия --- среднее число равновероятных вариантов следующего слова, между которыми модель как будто выбирает. Сравнивать перплексию корректно только на одинаковых тестовых данных и словаре.

\begin{important}
    \item Языковая модель оценивает вероятность текста или следующего слова.
    \item N-граммы ограничивают контекст последними $N-1$ словами.
    \item Главная проблема n-грамм --- разреженность и нулевые вероятности.
    \item Сглаживание делает вероятности ненулевыми и более устойчивыми.
    \item Перплексия ниже --- модель лучше предсказывает тестовый текст.
\end{important}
